\chapter{Welfare and the terms of sale}\label{ch:welfare}

\begin{chapterplan}
The question of this chapter: who gains when the challenger can see the price, and how do the payment rule and the
reserve change what the price does?

After this chapter, you can do these things:
\begin{itemize}
\item Define expected target proceeds and net acquisition surplus, and derive both with and without price access.
\item State and prove the proposition on access to prices, and name the hypotheses of Proposition~2 that it uses.
\item Derive the target's payment under Nash bargaining and the spread under every seller weight.
\item Say when a stronger incumbent widens the spread, narrows it, or leaves it unchanged under the payment rule.
\item State and prove the proposition on reserves, and explain why a reserve that excludes the low class makes
entry flat.
\end{itemize}
\end{chapterplan}

Chapter~\ref{ch:ccc} compared economies with different incumbents. This chapter keeps the incumbent fixed and changes
the terms of the sale. It covers three results of Section~7 of the paper. Their statements and proofs are in
Appendix~A.7 of the paper. They are the proposition on access to prices, the proposition on the payment rule, and the
proposition on reserves.

% ====================================================================
\section{The story in words}\label{sec:welfare-story}

Three groups take part in the sale of the target:
\begin{itemize}
\item the target's shareholders, who receive the terminal payoff \(V_T\) of their shares;
\item the two bidders, the incumbent and the challenger;
\item the traders in the order slot, the informed investor and the liquidity traders.
\end{itemize}
The first result asks who gains when the challenger sees the price. The second and third results ask how the payment
rule and the reserve change what the price can say.

\subsection*{Access to prices}

Compare two economies with the same strong incumbent. In the first economy, the challenger sees the price, as in
Chapter~\ref{ch:ccc}. In the second economy, it sees neither the price nor the order. The investor, the market makers
and the liquidity traders play the same game in both economies. Only the challenger's information differs.

In the first economy, a purchase is news. The price after a purchase tells the challenger that its class is probably
high. At this belief, gross profit covers the high cost, so a high-cost challenger prepares. In the second economy, the
challenger keeps the prior belief. Gross profit at the prior is below the high cost, so only the zero-cost challenger
prepares.

Access therefore adds high-cost preparation after a purchase, and nowhere else. These are the right challengers: each
added challenger expects a gross profit above its cost. Three groups gain.
\begin{itemize}
\item \emph{The target's shareholders.} A prepared challenger competes with the incumbent. The target receives more
on average, and it can receive a sale where it would otherwise stay unsold.
\item \emph{The economy as a whole.} An added preparation costs \(c_H\). It adds the value of a better match, or of
a sale that the reserve would otherwise block. The addition exceeds the cost.
\item \emph{The acquirers.} The added challenger expects a profit net of its cost. The incumbent loses part of its
profit, because a rival now competes for the same target. Net acquisition surplus already counts this loss, since it
measures the value of the final owner and not the profit of each bidder.
\end{itemize}
Section~\ref{sec:welfare-access} makes each statement exact.

\begin{analogy}[a queue at a ticket office]
A traveller pays to prepare a trip: it buys a guidebook, books leave, and plans the route. The trip is worth the cost
only if the destination is good. The traveller cannot see the destination, but it can see the queue at the ticket
office. Experienced travellers queue when the destination is good. A long queue is news, and the traveller prepares.

The analogy breaks in two places. Casual travellers join the queue at random, as liquidity traders do, so the queue is
a noisy sign. Also, an experienced traveller queues for its own profit and does not intend to inform anyone.
\end{analogy}

\subsection*{The payment rule and the reserve}

The seller announces the sale rule before trade starts. The rule decides two things. It divides the challenger's value
between the target and the winning bidder. It also decides how the spread \(\DT\) responds to competition, and the
spread is what the informed investor trades on. A seller who picks a payment rule or a reserve also picks how much
the price can say.

The payment rule changes the direction of the response. Under the second-price rule, a stronger incumbent widens the
spread. If the seller keeps a large share of the winner's value, a stronger incumbent narrows the spread.

The reserve changes the size of the response. A reserve that excludes the low class adds a term to the spread, and a
stronger incumbent shrinks this term. The weak incumbent then already leaves a wide spread. Competition has no price
left to switch on, and entry stays flat.

\begin{keyidea}
Price access changes only which high-cost challengers prepare. The orders, the trading costs and the zero-cost
challengers are the same with and without access.
\end{keyidea}

% ====================================================================
\section{What you need from the earlier chapters}\label{sec:welfare-setup}

This section recalls the objects of the chapter. Chapter~\ref{ch:model} gives the model, Chapter~\ref{ch:auction}
the auction, and Chapters~\ref{ch:beliefs} to~\ref{ch:ccc} the price and the orders.

\begin{itemize}
\item The incumbent's value \(R\) has a continuous distribution \(F\). The strong incumbent is \(F_1\), the weak
incumbent is \(F_0\), and \(F_1\le F_0\). Write \(\Delta_i=\DT(F_i)\) and \(B_i\) for the objects at \(F_i\).
\item The challenger is in class \(H\) with the prior probability \(\pi\), and in class \(L\) otherwise. Its cost
\(C\) is zero with probability \(\rho\) and equals \(c_H\) otherwise. Write \(\pi_H=\pi\) and \(\pi_L=1-\pi\).
\item Without entry the target receives \(t_0\) on average. After entry by class \(j\) it receives \(t_j\) on
average. Write \(w_j=t_j-t_0\) and \(w(\mu)=(1-\mu)w_L+\mu w_H=w_L+\DT\mu\). Then \(w_H-w_L=\DT\).
\item Gross profit at the belief \(\mu\) is \(B(\mu)=(1-\mu)g_L+\mu g_H\).
\item The slot gives a liquidity order in one direction with probability \(b=(1-\lambda)\beta\). The mass ratio is
\(u=\lambda/b\). The posteriors after a sale, no order and a purchase are \(m\), \(\pi\) and \(M\), and
\(\kappa=\min\{m,1-M\}\).
\item Entry after the order \(x\) is \(e(x)=\rho+(1-\rho)\zeta\), where \(\zeta\) is the probability that a high-cost
challenger prepares. Entry in class \(j\) is \(e_j=\sum_xf_j(x)\,e(x)\), where \(f_j(x)\) is the probability of the
order \(x\) in class \(j\).
\end{itemize}

Two formulas from Chapter~\ref{ch:beliefs} do most of the work. The terminal payoff of a share, given the class and
the order, is
\begin{equation}\label{eq:welfare-terminal}
\E[V_T\mid j,X=x]=t_0+e(x)\,w_j,
\end{equation}
which is \eqref{eq:beliefs-terminal}. The competitive price is \(P(x)=t_0+e(x)\,w(\mu(x))\), which is
\eqref{eq:beliefs-price}. Both hold in any economy, whatever the challenger sees. They need only truthful bids and a
cost that is independent of the class.

Chapter~\ref{ch:ccc} solved the economy at \(F_1\) under (H1) to (H3). I call it the access economy. The investor in
class \(H\) buys, and the investor in class \(L\) sells. The orders \((-1,0,1)\) have the class-\(H\) probabilities
\(f_H=(b,\,b_0,\,b+\lambda)\) and the class-\(L\) probabilities \(f_L=(b+\lambda,\,b_0,\,b)\). The posteriors are
\((m,\pi,M)\) and entry is \((\rho,\rho,1)\). Class entries are
\begin{equation}\label{eq:welfare-classentry}
e_H=\rho+(1-\rho)(b+\lambda),\qquad e_L=\rho+(1-\rho)\,b ,
\end{equation}
from Step~4 of the proof in Section~\ref{sec:ccc-proof-ii}.

% ====================================================================
\section{Access to prices}\label{sec:welfare-access}

\subsection{The two economies}\label{sec:welfare-twoeconomies}

The \emph{hidden-price economy} changes one thing. The challenger sees neither the price nor the order. It still
learns its cost. Market makers still see the order and still set competitive prices. So the price still moves with the
order, but the challenger cannot read it.

The challenger's belief is therefore the prior \(\pi\) after every order. Entry depends only on the cost, as
Table~\ref{tab:welfare-entry} shows.

\begin{table}[ht]
\centering
\begin{tabular}{@{}lccc@{}}
\toprule
Order \(x\) & sale, \(-1\) & no order, \(0\) & purchase, \(1\)\\
\midrule
Probability in class \(H\), \(f_H(x)\) & \(b\) & \(b_0\) & \(b+\lambda\)\\
Probability in class \(L\), \(f_L(x)\) & \(b+\lambda\) & \(b_0\) & \(b\)\\
Posterior \(\mu(x)\) & \(m\) & \(\pi\) & \(M\)\\
Entry with access & \(\rho\) & \(\rho\) & \(1\)\\
Entry without access & \(\rho\) & \(\rho\) & \(\rho\)\\
\bottomrule
\end{tabular}
\caption{The two economies at the strong incumbent \(F_1\), order by order. The orders and posteriors are the same.
Entry differs after a purchase only.}
\label{tab:welfare-entry}
\end{table}

The table claims that the hidden-price economy has these full orders and this entry in every equilibrium. The first
step of the proof shows it.

\subsection{Expected proceeds under any entry rule}\label{sec:welfare-proceeds}

The paper writes \(\mathcal R\) for expected target proceeds. This is the expected terminal payoff \(\E[V_T]\) of a
share. The orders and the entry rule fix it, as the next derivation shows. Average \eqref{eq:welfare-terminal} first
over the order, and then over the class.

\begin{derivation}[expected proceeds]
\begin{align*}
\mathcal R&=\E[V_T] \why{definition}\\
&=\sum_j\pi_j\sum_xf_j(x)\,\E[V_T\mid j,X=x] \why{total expectation over the class and the order}\\
&=\sum_j\pi_j\sum_xf_j(x)\,\bigl[t_0+e(x)\,w_j\bigr] \why{\eqref{eq:welfare-terminal}}\\
&=t_0\sum_j\pi_j\sum_xf_j(x)+\sum_j\pi_j\,w_j\sum_xf_j(x)\,e(x) \why{split the sum}\\
&=t_0+\sum_j\pi_j\,w_j\,e_j \why{\(\sum_xf_j(x)=1\), \(\sum_j\pi_j=1\), definition of \(e_j\)}
\end{align*}
\end{derivation}

So, for any orders and any entry rule,
\begin{equation}\label{eq:welfare-R}
\mathcal R=t_0+\pi\,w_H\,e_H+(1-\pi)\,w_L\,e_L .
\end{equation}
In words: the target receives \(t_0\), plus the gain \(w_j\) from entry, weighted by the chance that the challenger is
in class \(j\) and prepares.

\subsection{Net acquisition surplus}\label{sec:welfare-surplus}

The paper measures welfare by acquisition surplus net of preparation costs, \(\mathcal W\). Surplus is the value of the
buyer that ends with the target. A payment moves money from a buyer to the shareholders, so it does not enter. A
trading profit moves money between traders, so it does not enter either. Preparation costs are real resources, so
they enter with a minus sign.

Define the allocation value of a sale. If no bidder is admissible, the target stays unsold and the value is zero. If
one or more bidders are admissible, the value is the highest admissible value. Under the class ranking, every prepared
challenger is admissible, since \(\theta\ge p\). Let
\[
\mathcal V_0(R)=R\,\1\{R\ge p\}
\]
be the allocation value without a challenger. With a prepared challenger of value \(\theta\), it is
\(\max\{\mathcal V_0(R),\theta\}\), because bids are truthful and the target goes to the highest value.

\begin{paperresult}{The gain from one prepared challenger (the paper's identity, Appendix A.7)}
For \(\theta\ge p\) and every \(R\ge0\),
\begin{equation}\label{eq:welfare-identity}
\max\{\mathcal V_0(R),\theta\}-\mathcal V_0(R)=(\theta-\max\{p,R\})_++p\,\1\{R<p\}.
\end{equation}
\end{paperresult}

The positive part \((x)_+=\max\{x,0\}\) is the paper's notation. The proof is a check of two cases.

\begin{derivation}[the identity \eqref{eq:welfare-identity}]
\emph{Case \(R<p\).} The incumbent is not admissible, so \(\mathcal V_0(R)=0\).
\begin{align*}
\text{left side}&=\max\{0,\theta\}-0 \why{\(\mathcal V_0(R)=0\)}\\
&=\theta \why{\(\theta\ge p>0\)}\\
\text{right side}&=(\theta-p)_++p \why{\(\max\{p,R\}=p\), \(\1\{R<p\}=1\)}\\
&=(\theta-p)+p=\theta \why{\(\theta\ge p\)}
\end{align*}
\emph{Case \(R\ge p\).} The incumbent is admissible, so \(\mathcal V_0(R)=R\).
\begin{align*}
\text{left side}&=\max\{R,\theta\}-R \why{\(\mathcal V_0(R)=R\)}\\
&=(\theta-R)_+ \why{\(\max\{R,\theta\}-R=\max\{0,\theta-R\}\)}\\
\text{right side}&=(\theta-R)_++0 \why{\(\max\{p,R\}=R\), \(\1\{R<p\}=0\)}
\end{align*}
\end{derivation}

Average \eqref{eq:welfare-identity} over \(R\). The first term has the mean \(g(\theta)\), the challenger's expected
gross profit (Section~\ref{sec:toolkit-tail} gives the tail formula). The second term has the mean \(pF(p)\).
\begin{align*}
\E_R\bigl[(\theta-\max\{p,R\})_+\bigr]&=g(\theta) \why{definition of \(g\)}\\
\E_R\bigl[p\,\1\{R<p\}\bigr]&=p\,F(p) \why{\(\Pr(R<p)=F(p)\), \(F\) continuous}
\end{align*}
Average over the class as well. Under the class ranking, every class-\(j\) value meets the reserve. The mean gain from
one prepared class-\(j\) challenger is therefore
\begin{equation}\label{eq:welfare-aj}
a_j=g_j+p\,F(p).
\end{equation}
The first part, \(g_j\), is the challenger's own gross profit. The second part, \(pF(p)\), is a sale that exists only
because the challenger came: it happens when the incumbent fails the reserve, which has probability \(F(p)\), and the
challenger buys at \(p\).

Now add up over the classes and the orders. A high-cost challenger that prepares pays \(c_H\). The probability that a
class-\(j\) challenger prepares with the high cost is \(e_j-\rho\), because entry \(e_j\) includes the zero-cost
challengers, who have probability \(\rho\) and pay nothing.

\begin{derivation}[net acquisition surplus under any entry rule]
Let \(v_0=\E[\mathcal V_0(R)]\). Given the class and the order, the challenger prepares with probability \(e(x)\),
and preparation is independent of \((\theta,R)\). So
\begin{align*}
\mathcal W&=\sum_j\pi_j\sum_xf_j(x)\,\bigl[v_0+e(x)\,a_j\bigr]-c_H\sum_j\pi_j\,(e_j-\rho)
\why{allocation value minus expected preparation cost}\\
&=v_0+\sum_j\pi_j\,a_j\,e_j-c_H\sum_j\pi_j\,(e_j-\rho) \why{\(\sum_xf_j(x)=1\), definition of \(e_j\)}\\
&=v_0+\sum_j\pi_j\bigl\{a_j\,e_j-c_H\,(e_j-\rho)\bigr\} \why{combine the sums}
\end{align*}
\end{derivation}

This is the surplus counterpart of \eqref{eq:welfare-R}:
\begin{equation}\label{eq:welfare-W}
\mathcal W=v_0+\sum_j\pi_j\bigl\{a_j\,e_j-c_H\,(e_j-\rho)\bigr\}.
\end{equation}
Each formula has the same shape. A part of the economy does not depend on entry. The rest is the entry in each class
times the gain from one more prepared challenger of that class.

\subsection{The proposition}\label{sec:welfare-accessstatement}

\needspace{12\baselineskip}
\begin{paperresult}{The proposition on access to prices (Appendix A.7 of the paper)}
Hold the incumbent at \(F_1\). Let the hypotheses of Proposition~2 hold. Compare the equilibrium of Proposition~2 with
the equilibrium of an economy in which the challenger does not see the price. All other primitives stay fixed, and
market makers set competitive prices in both economies. Write \(\mathcal R\) for expected target proceeds and
\(\mathcal W\) for expected acquisition surplus net of preparation costs. The superscript \(1\) marks price access
and \(0\) marks the hidden price. Acquisition surplus excludes transfers between bidders, shareholders and traders.

Then both economies have unique full orders. Price access raises expected target proceeds and acquisition surplus net
of preparation costs:
\[
\begin{aligned}
\mathcal R^1-\mathcal R^0&=(1-\rho)\bigl\{\pi(b+\lambda)(t_H-t_0)+(1-\pi)\,b\,(t_L-t_0)\bigr\}>0,\\
\mathcal W^1-\mathcal W^0&=(1-\rho)(b+\pi\lambda)\bigl[B_1(M)-c_H+pF_1(p)\bigr]>0.
\end{aligned}
\]
\end{paperresult}

Read the statement line by line.
\begin{itemize}
\item \emph{The hypotheses of Proposition~2.} They are the class ranking with finite means, continuous
\(F_2\le F_1\le F_0\), interior \(\pi,\rho,\lambda\), \(0<\beta<1/2\), \(c_H>0\), \(k>0\), and (H1) to (H3). The proof
below uses the class ranking, \(F_1\le F_0\) and continuity, the two sides of (H1), and the right side of (H2),
\(k<\rho\kappa\Delta_1\). It does not use \(\Delta_0<k\), (H3) or \(F_2\). The incumbent \(F_0\) enters only through
\(B_1(\pi)\le B_0(\pi)<c_H\). This list is a reading of the proof by the guide, not a statement of the paper.
\item \emph{Unique full orders.} In each economy, class \(H\) buys and class \(L\) sells, in every equilibrium.
\item \emph{The first line.} The gain in proceeds is the chance \(1-\rho\) of a high cost, times the chance of a
purchase in each class, times the entry gain \(w_j\) of that class.
\item \emph{The second line.} The gain in surplus is the chance of an added preparation, \((1-\rho)(b+\pi\lambda)\),
times the net gain of one added preparation.
\end{itemize}

\subsection{The proof}\label{sec:welfare-accessproof}

\begin{derivation}[proof of the proposition on access to prices]
\begin{enumerate}[label=\textit{Step \arabic*.},leftmargin=*,itemsep=0.6em]
\item \emph{The access economy.} It is the economy of Proposition~2 at \(F_1\). The proposition gives a unique
equilibrium with full orders, posteriors \((m,\pi,M)\) and entry \((\rho,\rho,1)\). The class entries are
\eqref{eq:welfare-classentry}.

\item \emph{Every candidate of the hidden-price economy has entry \(\rho\).} A candidate meets Bayes' rule for the
posteriors, competitive prices, and truthful bids. The challenger's belief equals the prior \(\pi\), since it learns
nothing from the price or the order. By Proposition~1(a), a stronger incumbent never raises gross profit at a fixed
belief. With (H1), this gives
\begin{align*}
B_1(\pi)&\le B_0(\pi) \why{Proposition~1(a), \(F_1\le F_0\)}\\
&<c_H \why{left side of (H1)}
\end{align*}
Gross profit at the challenger's belief is below the high cost, strictly. A high-cost challenger does not prepare, and
no tie arises. So \(e(x)=\rho\) after every order, in every candidate.

\item \emph{Every candidate of the hidden-price economy has full orders.} Use the trading returns
\eqref{eq:trading-returns} at \(F_1\), with \(e\equiv\rho\):
\begin{align*}
K_H&=\rho\,\Delta_1\,[1-\mu(1)] \why{\eqref{eq:trading-returns}, \(e(1)=\rho\)}\\
&\ge\rho\,\Delta_1\,(1-M) \why{\(\mu(1)\le M\) under every strategy, \(\Delta_1>0\)}\\
&\ge\rho\,\kappa\,\Delta_1 \why{\(1-M\ge\kappa\)}\\
&>k \why{right side of (H2)}
\end{align*}
In the same way \(K_L=\rho\Delta_1\mu(-1)\ge\rho\Delta_1m\ge\rho\kappa\Delta_1>k\). Here \(\Delta_1>0\) because
\(\Delta_1>k/(\rho\kappa)>0\). A wrong-signed order earns at most \(-k\). So the correct order is the unique best
order of each class, and \(\sigma_H=\delta_1\), \(\sigma_L=\delta_{-1}\). The bounds hold for every candidate and need
no tie rule.

\item \emph{The hidden-price tuple.} Steps~2 and~3 fix a unique candidate. Its posteriors are \((m,\pi,M)\), its
entry is \(\rho\), its belief is \(\pi\), and its prices are \(P^0(x)=t_0+\rho\,w_1(\mu(x))\). It meets every
condition: Bayes' rule holds by the order probabilities, the belief is the prior by construction, the high-cost
choice is strict by Step~2, the prices are competitive by the price formula with \(e\equiv\rho\), and the orders are
strict best orders by Step~3. So the hidden-price economy has exactly one equilibrium, and its orders are full.

\item \emph{The orders coincide.} By Steps~1 and~4, both economies have the probabilities \(f_H\) and \(f_L\) of
Table~\ref{tab:welfare-entry}. The trading costs coincide. Entry differs after a purchase only.

\item \emph{The proceeds gain.} Entry in the hidden-price economy is \(e_H^0=e_L^0=\rho\). Entry in the access
economy is \eqref{eq:welfare-classentry}. Apply \eqref{eq:welfare-R} in both:
\begin{align*}
\mathcal R^1-\mathcal R^0&=\pi\,w_H\,(e_H^1-\rho)+(1-\pi)\,w_L\,(e_L^1-\rho) \why{\eqref{eq:welfare-R} twice; \(t_0\) cancels}\\
&=\pi\,w_H\,(1-\rho)(b+\lambda)+(1-\pi)\,w_L\,(1-\rho)\,b \why{\eqref{eq:welfare-classentry}}\\
&=(1-\rho)\bigl\{\pi(b+\lambda)(t_H-t_0)+(1-\pi)\,b\,(t_L-t_0)\bigr\} \why{\(w_j=t_j-t_0\), factor \(1-\rho\)}
\end{align*}

\item \emph{The surplus gain.} Apply \eqref{eq:welfare-W} in both economies. The terms \(v_0\) and \(c_H\rho\) cancel:
\begin{align*}
\mathcal W^1-\mathcal W^0&=\sum_j\pi_j\,(a_j-c_H)\,(e_j^1-\rho) \why{\eqref{eq:welfare-W} twice; collect \(e_j\)}\\
&=(1-\rho)\bigl\{\pi(b+\lambda)(a_H-c_H)+(1-\pi)\,b\,(a_L-c_H)\bigr\} \why{\eqref{eq:welfare-classentry}}
\end{align*}
Convert the class weights to the posterior \(M\). A purchase has probability
\(\Pr(X=1)=\pi(b+\lambda)+(1-\pi)\,b=b+\pi\lambda\). Bayes' rule gives \(M=\pi(b+\lambda)/\Pr(X=1)\), so
\begin{align*}
\pi(b+\lambda)&=(b+\pi\lambda)\,M \why{definition of \(M\)}\\
(1-\pi)\,b&=(b+\pi\lambda)\,(1-M) \why{complement}
\end{align*}
Substitute, and use \eqref{eq:welfare-aj}:
\begin{align*}
\mathcal W^1-\mathcal W^0&=(1-\rho)(b+\pi\lambda)\bigl\{M(a_H-c_H)+(1-M)(a_L-c_H)\bigr\} \why{substitute the weights}\\
&=(1-\rho)(b+\pi\lambda)\bigl\{M\,a_H+(1-M)\,a_L-c_H\bigr\} \why{\(M+(1-M)=1\)}\\
&=(1-\rho)(b+\pi\lambda)\bigl\{B_1(M)+pF_1(p)-c_H\bigr\} \why{\eqref{eq:welfare-aj}, definition of \(B_1\)}
\end{align*}
The same conversion gives the proceeds gain as \((1-\rho)(b+\pi\lambda)\,w_1(M)\).

\item \emph{The signs.} The factors \(1-\rho\) and \(b+\pi\lambda\) are positive. For proceeds, \(w_L\ge0\) because
entry never lowers proceeds (Section~\ref{sec:auction-outcomes}), and
\begin{align*}
w_H&=w_L+\Delta_1 \why{\(w_H-w_L=\DT\)}\\
&\ge\Delta_1>0 \why{\(w_L\ge0\), Step~3}
\end{align*}
So the term \(\pi(b+\lambda)w_H\) is positive and the other term is nonnegative. For surplus,
\(B_1(M)-c_H>0\) by the right side of (H1), and \(pF_1(p)\ge0\).
\end{enumerate}
\end{derivation}

\paragraph{What this says.} Access works only through the chance that a high-cost challenger prepares after a
purchase. That chance is \((1-\rho)\Pr(X=1)\). Each such preparation is worth \(w_1(M)\) to the target and
\(B_1(M)+pF_1(p)-c_H\) to the economy.

\paragraph{Extreme cases.}
\begin{itemize}
\item If \(\rho\) is near one, almost every challenger has zero cost, and both gains tend to zero. Nothing is left
for the price to recruit.
\item If the incumbent never meets the reserve, then \(F_1(p)=1\). A challenger always buys at \(p\), and the whole
surplus gain from an added preparation is its value minus \(c_H\). The reserve term \(pF_1(p)\) is at its largest.
\item If gross profit at \(M\) only just exceeds \(c_H\), the term \(B_1(M)-c_H\) is near zero. The surplus gain is
then the sale term, which is nonnegative. Access cannot lower surplus here.
\end{itemize}

\begin{watchout}
The proposition compares the two economies at \(F_1\), and in each economy the equilibrium is solved again. A
deviation in one economy faces the prices and entry of that economy. The proposition does not compare incumbents.
It also does not say that every preparation helps: it says that preparation after a purchase does, because
\(B_1(M)>c_H\). The next paragraph shows the other case.
\end{watchout}

\subsection{Who gets the surplus gain}\label{sec:welfare-split}

The paper states the gain in proceeds and the gain in surplus. It does not state the gain of each bidder. The guide
adds the following accounting, which uses only the paper's formulas.

The value of the final owner equals the price it pays plus its profit. The price is what the target receives. So, in
every sale, allocation value equals target proceeds plus the incumbent's profit plus the challenger's profit, where
the loser's profit is zero. Take one added preparation, after a purchase, so that the class is \(H\) with
probability \(M\). The three parts change as follows.
\begin{itemize}
\item Target proceeds change by \(w_1(M)\), the average of \(w_H\) and \(w_L\) at the weights \(M\) and \(1-M\).
\item The challenger's profit before its cost is \(B_1(M)\). Net of the cost, it is \(B_1(M)-c_H\), which is positive.
\item The incumbent's profit changes by the remainder.
\end{itemize}
The sum must equal the gain in allocation value, \(B_1(M)+pF_1(p)\). So the remainder is
\begin{align*}
\Delta(\text{incumbent profit})&=\bigl[B_1(M)+pF_1(p)\bigr]-w_1(M)-B_1(M) \why{allocation value minus the other two parts}\\
&=-\bigl[w_1(M)-pF_1(p)\bigr] \why{\(B_1(M)\) cancels}
\end{align*}
This loss is not positive. For \(\theta\ge p\), the proceeds after entry are \(t(\theta)=p+\int_p^\theta[1-F_1(y)]\,dy\ge p\)
(Section~\ref{sec:auction-spread}). So \(t_j\ge p\), and
\begin{align*}
w_j&=t_j-t_0 \why{definition}\\
&\ge p-p\,[1-F_1(p)] \why{\(t_j\ge p\), \(t_0=p[1-F_1(p)]\)}\\
&=p\,F_1(p) \why{expand}
\end{align*}
Both classes satisfy it, so \(w_1(M)\ge pF_1(p)\). The incumbent loses at most the part of the proceeds gain that
the sale term does not explain.

\subsection{The investor's profit and the shareholders}\label{sec:welfare-transfer}

The informed investor earns a trading profit in the access economy. This profit is a transfer from the liquidity
traders. The paper computes it, and it asks whether the initial shareholders still gain if they bear all of it.

Class \(H\) buys and class \(L\) sells. Before the cost \(k\), a correct order earns the trading return
\eqref{eq:trading-returns}. The investor arrives with probability \(\lambda\). Its expected gross profit is
\begin{align*}
G&=\lambda\bigl[\pi\,K_H+(1-\pi)\,K_L\bigr] \why{arrival, class weights}\\
&=\lambda\Delta_1\bigl\{\pi\,(1-M)+(1-\pi)\,\rho\,m\bigr\} \why{\(K_H=1\cdot\Delta_1(1-M)\), \(K_L=\rho\Delta_1m\)}
\end{align*}
The returns use entry \(e(1)=1\) after a purchase and \(e(-1)=\rho\) after a sale. We have \(G>0\).

\emph{The liquidity traders lose \(G\).} Competitive prices give \(\E[V_T-P(X)\mid X]=0\). Hence
\begin{align*}
\E\bigl[X\,(V_T-P(X))\bigr]&=\sum_xx\,\Pr(X=x)\,\E[V_T-P(x)\mid X=x] \why{total expectation over the order}\\
&=0 \why{\(P(x)=\E[V_T\mid X=x]\)}
\end{align*}
The order \(X\) comes from the informed investor or from a liquidity trader, so the expectation splits into two
parts. The informed part is \(G\). So the liquidity part is \(-G\).

Now measure one order unit as \(n\) shares per share of float. The economy without an informed investor has the
posterior \(\pi\) after every order, so entry is \(\rho\) and its proceeds equal \(\mathcal R^0\). If the initial
shareholders bear the whole transfer, their total gain relative to that economy is
\[
\mathcal R^1-\mathcal R^0-n\,G .
\]
It is positive exactly when \(n<(\mathcal R^1-\mathcal R^0)/G\). The ratio depends only on the primitives, and
Chapter~\ref{app:numbers} gives its value in the example.

\begin{watchout}
The break-even order unit rests on an extreme assumption: the initial shareholders bear the whole transfer. In the
model, a liquidity sale comes from a current shareholder, and a liquidity purchase comes from an outside investor. So
the assumption is conservative for the shareholders. The result says that they still gain when the order unit is below
the ratio. It does not say how the transfer is shared.
\end{watchout}

\subsection{Information, not the price level}\label{sec:welfare-dividend}

Access raises the mean price. By competitive pricing, \(\E[P(X)]=\E[V_T]=\mathcal R\), so the mean price in the access
economy exceeds the mean price in the hidden-price economy by \(\mathcal R^1-\mathcal R^0\). One could suspect that the
higher price level draws entry, and not the information in the price. The paper answers with a control.

Add a deterministic external dividend \(d\) to every share in the hidden-price economy, and take
\(d=\mathcal R^1-\mathcal R^0\). Check each object in turn.
\begin{align*}
V_T&\to V_T+d \why{the dividend}\\
P(x)=\E[V_T\mid X=x]&\to P(x)+d \why{\(d\) is a constant}\\
q\,[V_T-P(q)]-k|q|&\to q\,[(V_T+d)-(P(q)+d)]-k|q| \why{shift both terms}\\
&=q\,[V_T-P(q)]-k|q| \why{\(d\) cancels}
\end{align*}
Every order payoff is unchanged, so the orders are unchanged. The challenger sees nothing, and the dividend is no part of
the sale, so entry is unchanged. Proceeds of the sale are unchanged as well. The dividend raises the mean price to
\(\mathcal R^0+d=\mathcal R^1\). The two economies now have the same mean price, and they still differ in entry and in
proceeds. They differ in the information in the price, and not in its level.

With costly research (Chapter~\ref{ch:research}), the paper reports the same gains in Appendix A.5. This chapter does
not derive that case.

% ====================================================================
\subsection{Access outside the recruitment region}\label{sec:welfare-outside}

Proposition~2 puts the economy inside the recruitment region: gross profit is below \(c_H\) at the prior and above
\(c_H\) at the top posterior. The paper also looks at the opposite case. Access then removes preparation.

Suppose that
\[
B_1(m)<c_H<B_1(\pi)\qquad\text{and}\qquad k<\rho\kappa\Delta_1 .
\]
Gross profit covers the high cost at the prior but not after a sale.
\begin{itemize}
\item \emph{Without the price.} The challenger's belief is \(\pi\). Gross profit \(B_1(\pi)\) exceeds \(c_H\), so every
challenger prepares: entry is one after every order.
\item \emph{With the price.} The posteriors are \((m,\pi,M)\). Gross profit is nondecreasing in the belief, since
\(g_H\ge g_L\). So \(B_1(M)\ge B_1(\pi)>c_H\) and \(B_1(m)<c_H\). Entry is \((\rho,1,1)\).
\item \emph{The orders.} In both economies entry lies in \([\rho,1]\), so the argument of Step~3 of the proof applies
and the orders are full. The order probabilities of Table~\ref{tab:welfare-entry} hold in both.
\end{itemize}
The paper states that each tuple then meets all six equilibrium conditions, as in the proof above. Access removes
high-cost preparation after a sale, and nowhere else. A sale has probability \(b\) in class \(H\) and \(b+\lambda\) in
class \(L\).

\begin{derivation}[the proceeds effect of access outside the recruitment region]
With access, class entries add the entry after each order, weighted by the order probabilities. The sum of the three
class-\(H\) probabilities is one, \(b+b_0+(b+\lambda)=1\), and the class-\(L\) probabilities also sum to one.
\begin{align*}
e_H^1&=\rho\,b+1\cdot b_0+1\cdot(b+\lambda) \why{entry \((\rho,1,1)\), class-\(H\) probabilities}\\
&=1-(1-\rho)\,b \why{\(b+b_0+(b+\lambda)=1\)}\\
e_L^1&=\rho\,(b+\lambda)+b_0+b=1-(1-\rho)(b+\lambda) \why{same steps, class-\(L\) probabilities}
\end{align*}
Entry without access is one in both classes. Apply \eqref{eq:welfare-R} in both economies:
\begin{align*}
\mathcal R^1-\mathcal R^0&=\pi\,w_H\,(e_H^1-1)+(1-\pi)\,w_L\,(e_L^1-1) \why{\eqref{eq:welfare-R} twice; \(t_0\) cancels}\\
&=-(1-\rho)\bigl\{\pi\,b\,w_H+(1-\pi)(b+\lambda)\,w_L\bigr\} \why{the entry changes}
\end{align*}
\end{derivation}

The proceeds fall. The first term in the braces is positive, since \(\pi b\,w_H>0\), and the second is nonnegative. The
same steps as in Step~6 of the proof, applied to \eqref{eq:welfare-W} and to a sale in place of a purchase, give the
surplus change
\[
\mathcal W^1-\mathcal W^0=-(1-\rho)\bigl[b+(1-\pi)\lambda\bigr]\bigl[B_1(m)-c_H+pF_1(p)\bigr].
\]
Question~\ref{q:welfare-sign} asks you to fill in the steps. The paper reports the numbers for the example with the high
cost raised to the level that puts the economy in this case.

\begin{checkyourself}\label{q:welfare-sign}
Derive the surplus change above from \eqref{eq:welfare-W}. Then say when access raises surplus and when it lowers
surplus. Which part of the expression can make a removed preparation worth its cost to the economy, although it is not
worth its cost to the challenger?
\end{checkyourself}

% ====================================================================
\section{The payment rule}\label{sec:welfare-payment}

\subsection{The family of rules}\label{sec:welfare-family}

The second-price auction pays the target the lower of the two values. Other sale institutions pay it more. The paper
studies a family indexed by one number, the seller weight \(\omega\in[0,1)\).

The setting is a \emph{verifiable-value institution with zero reserve}. After diligence, the values of the two buyers
are verifiable. A sale to the next buyer, at its value, is an enforceable fallback. The seller and the winner divide the
surplus above the fallback by Nash bargaining, with the weight \(\omega\) on the seller. This is not the cash
second-price auction of the main model. The proposition studies only the payment stage: given the two values, what does
the target receive, and what does the challenger keep? It does not solve trade or entry, so no tie rule enters.

Write \(a\) for the lower of the two values and \(c\) for the higher. Then \(a\) is the fallback and \(c-a\) is the
surplus above it. The rule gives the target \(a+\omega(c-a)\). Three cases show the range.
\begin{itemize}
\item \(\omega=0\): the target receives \(a\). This is the second-price rule with a zero reserve.
\item \(\omega=\tfrac12\): the seller and the winner split the surplus equally.
\item \(\omega\) near one: the target receives almost \(c\), the winner's value.
\end{itemize}
Without a challenger the fallback is zero, and the target receives \(\omega R\).

\subsection{The transfer under Nash bargaining}\label{sec:welfare-nash}

Fix the incumbent's value \(R\) and the challenger's value \(\theta\). Write \(a=\min\{R,\theta\}\) and
\(c=\max\{R,\theta\}\). Take \(a<c\) and \(0<\omega<1\). The target receives a transfer \(T\) with \(a\le T\le c\). Its
gain over the fallback is \(T-a\). The winner's gain is \(c-T\). Nash bargaining maximizes the product of the two gains,
each raised to the weight of its owner:
\[
N(T)=\bigl(T-a\bigr)^{\omega}\,\bigl(c-T\bigr)^{1-\omega}.
\]

\begin{derivation}[the Nash transfer]
Maximize the logarithm of \(N\). The logarithm rises with \(N\), so both have the same maximizer.
\begin{align*}
\ln N(T)&=\omega\ln(T-a)+(1-\omega)\ln(c-T) \why{log of a product}\\
\frac{d\ln N}{dT}&=\frac{\omega}{T-a}-\frac{1-\omega}{c-T} \why{derivative of each logarithm}\\
\frac{d^2\ln N}{dT^2}&=-\frac{\omega}{(T-a)^2}-\frac{1-\omega}{(c-T)^2}<0 \why{derivative again; \(\omega,1-\omega>0\)}
\end{align*}
The second derivative is negative on \(a<T<c\), so there is at most one maximum, and it solves the first-order
condition.
\begin{align*}
\frac{\omega}{T-a}&=\frac{1-\omega}{c-T} \why{first derivative equals zero}\\
\omega\,(c-T)&=(1-\omega)\,(T-a) \why{cross-multiply}\\
\omega\,c+(1-\omega)\,a&=\omega\,T+(1-\omega)\,T \why{move the terms in \(T\) to one side}\\
\omega\,c+(1-\omega)\,a&=T \why{\(\omega+(1-\omega)=1\)}
\end{align*}
The log tends to minus infinity at both ends of the interval, so this interior point is the maximum.
\end{derivation}

Two edge cases need a word. If \(R=\theta\), then \(a=c\), and feasibility forces \(T=\theta\), which the formula also
gives. If \(\omega=0\), the product is \(c-T\), the next buyer accepts its fallback, and \(T=a\), which the formula also
gives. In every case,
\begin{equation}\label{eq:welfare-transfer}
T_\omega(R,\theta)=(1-\omega)\min\{R,\theta\}+\omega\max\{R,\theta\}.
\end{equation}
In words: the target receives the lower value with weight \(1-\omega\) and the higher value with weight \(\omega\).
That is another way to write \(a+\omega(c-a)\).

\subsection{The challenger's profit and the spread}\label{sec:welfare-spreadpay}

The challenger wins when \(\theta>R\). It then keeps its value minus the transfer. Use \eqref{eq:welfare-transfer} with
\(\min\{R,\theta\}=R\) and \(\max\{R,\theta\}=\theta\):
\begin{align*}
\theta-T_\omega(R,\theta)&=\theta-(1-\omega)R-\omega\theta \why{\eqref{eq:welfare-transfer}, \(\theta>R\)}\\
&=(1-\omega)(\theta-R) \why{\(\theta-\omega\theta=(1-\omega)\theta\)}
\end{align*}
When \(\theta<R\) the challenger loses and keeps nothing. So its expected profit, with the value \(\theta\), is
\begin{equation}\label{eq:welfare-profitpay}
g_{\theta,\omega}=(1-\omega)\,\E\bigl[(\theta-R)_+\bigr].
\end{equation}

Now take the point-mass classes of the example: class \(H\) has the value \(h\), and class \(L\) has the value \(\ell\).
Let the incumbent's value lie in \([0,\bar r]\) with \(0<\ell<\bar r<h\). The spread is the difference of the two
expected transfers, so find the difference at a fixed \(R\). Two cases cover all \(R\), since \(R<h\).

\begin{derivation}[the class difference at a fixed incumbent value]
\emph{Case \(R\le\ell\).} Both classes win or tie. For each class the lower value is \(R\).
\begin{align*}
T_\omega(R,h)-T_\omega(R,\ell)&=\bigl[(1-\omega)R+\omega h\bigr]-\bigl[(1-\omega)R+\omega\ell\bigr]
\why{\eqref{eq:welfare-transfer}, \(R\le\ell<h\)}\\
&=\omega\,(h-\ell) \why{\(R\) cancels}
\end{align*}
\emph{Case \(\ell<R<h\).} Class \(H\) wins, and the lower value is \(R\). Class \(L\) loses to the incumbent, and the
lower value is \(\ell\).
\begin{align*}
T_\omega(R,h)-T_\omega(R,\ell)&=\bigl[(1-\omega)R+\omega h\bigr]-\bigl[(1-\omega)\ell+\omega R\bigr]
\why{\eqref{eq:welfare-transfer}}\\
&=(1-2\omega)R+\omega h-(1-\omega)\ell \why{collect \(R\): \((1-\omega)-\omega=1-2\omega\)}\\
&=(1-2\omega)(R-\ell)+\omega h-(1-\omega)\ell+(1-2\omega)\ell \why{add and subtract \((1-2\omega)\ell\)}\\
&=(1-2\omega)(R-\ell)+\omega\,(h-\ell) \why{\(-(1-\omega)\ell+(1-2\omega)\ell=-\omega\ell\)}
\end{align*}
\end{derivation}

One formula covers both cases, since \((R-\ell)_+=0\) in the first case:
\[
T_\omega(R,h)-T_\omega(R,\ell)=\omega\,(h-\ell)+(1-2\omega)(R-\ell)_+ .
\]
Take the expectation over \(R\). The spread is the expected class difference:
\begin{equation}\label{eq:welfare-spreadpay}
\Delta_\omega=\omega\,(h-\ell)+(1-2\omega)\,\E\bigl[(R-\ell)_+\bigr].
\end{equation}

\begin{watchout}
The hypothesis \(\bar r<h\) removes a third case. If the incumbent could exceed \(h\), both classes would lose. The
difference is then \((1-\omega)h+\omega R-(1-\omega)\ell-\omega R=(1-\omega)(h-\ell)\). The formula would hold with
\(\min\{(R-\ell)_+,h-\ell\}\) in place of \((R-\ell)_+\). The proposition avoids this cap.
\end{watchout}

\subsection{The proposition}\label{sec:welfare-paystatement}

\needspace{12\baselineskip}
\begin{paperresult}{The proposition on the payment rule (Appendix A.7 of the paper)}
Consider a verifiable-value institution with zero reserve. After diligence, values are verifiable, and a sale to the
next buyer at its value is an enforceable fallback. The seller and the winner divide the surplus above the fallback by
Nash bargaining, with seller weight \(\omega\in[0,1)\). Without a challenger, the seller receives \(\omega R\).

Let the challenger's values be the point masses \(\ell\) and \(h\), and let the incumbent distribution have support
\([0,\bar r]\) with \(0<\ell<\bar r<h\). The transfer, the challenger's expected profit and the spread are
\[
\begin{aligned}
T_\omega(R,\theta)&=(1-\omega)\min\{R,\theta\}+\omega\max\{R,\theta\},\\
g_{\theta,\omega}&=(1-\omega)\,\E[(\theta-R)_+],\\
\Delta_\omega&=\omega(h-\ell)+(1-2\omega)\,\E[(R-\ell)_+].
\end{aligned}
\]
A stronger incumbent weakly lowers the challenger's profit. Its effect on the spread weakly has the sign of
\(1-2\omega\). The effect is strict when \(\E[(R-\ell)_+]\) changes strictly and \(\omega\ne1/2\). At \(\omega=1/2\),
the spread equals \((h-\ell)/2\) at every incumbent distribution.
\end{paperresult}

Read the statement line by line.
\begin{itemize}
\item \emph{The three formulas} are \eqref{eq:welfare-transfer}, \eqref{eq:welfare-profitpay} and
\eqref{eq:welfare-spreadpay}. They are derived above.
\item \emph{The challenger's profit.} A stronger incumbent never raises it, for every \(\omega\in[0,1)\).
\item \emph{The spread.} The direction of its change is the sign of \(1-2\omega\). For \(\omega<1/2\) the spread
weakly rises. For \(\omega>1/2\) it weakly falls.
\item \emph{Strictness.} A strict change needs two things: the positive-part expectation \(\E[(R-\ell)_+]\) changes
strictly, and \(\omega\ne1/2\).
\item \emph{The midpoint.} At \(\omega=1/2\), the coefficient \(1-2\omega\) is zero and the spread is constant.
\end{itemize}

The model's incumbent distribution is continuous. A stronger incumbent is a continuous \(\tilde F\le F\)
(Chapter~\ref{ch:auction}). Take \(\bar r<h\) large enough for both distributions, so that both equal one above \(\bar r\).

\begin{derivation}[proof of the proposition on the payment rule]
The three formulas are \eqref{eq:welfare-transfer}, \eqref{eq:welfare-profitpay} and \eqref{eq:welfare-spreadpay}.
It remains to prove the statements about a stronger incumbent.
\begin{enumerate}[label=\textit{Step \arabic*.},leftmargin=*,itemsep=0.6em]
\item \emph{The profit.} Fix \(\theta>0\). The positive part \((\theta-R)_+\) is the length of the levels \(y\in[0,\theta]\)
that lie above \(R\). This is the fixed-value tail formula (Section~\ref{sec:toolkit-tail}) with a zero reserve.
\begin{align*}
\E\bigl[(\theta-R)_+\bigr]&=\E\Bigl[\int_0^\theta\1\{R<y\}\,dy\Bigr] \why{length of the levels above \(R\)}\\
&=\int_0^\theta\Pr(R<y)\,dy \why{swap; the integrand is nonnegative}\\
&=\int_0^\theta F(y)\,dy \why{\(F\) is continuous}
\end{align*}
The same holds for \(\tilde F\). The change in profit is
\begin{align*}
g_{\theta,\omega}(F)-g_{\theta,\omega}(\tilde F)&=(1-\omega)\int_0^\theta\bigl[F(y)-\tilde F(y)\bigr]\,dy
\why{\eqref{eq:welfare-profitpay}, the formula twice}\\
&\ge0 \why{\(1-\omega>0\) and \(F\ge\tilde F\)}
\end{align*}
\item \emph{The positive part.} By \eqref{eq:toolkit-pospart}, \(\E[(R-\ell)_+]=\int_\ell^\infty[1-F(y)]\,dy\). Then
\begin{align*}
\E_{\tilde F}\bigl[(R-\ell)_+\bigr]-\E_F\bigl[(R-\ell)_+\bigr]&=\int_\ell^{\bar r}\bigl[F(y)-\tilde F(y)\bigr]\,dy
\why{the formula twice; combine; both integrands vanish above \(\bar r\)}\\
&\ge0 \why{\(F\ge\tilde F\)}
\end{align*}
\item \emph{The spread.} Subtract \eqref{eq:welfare-spreadpay} at \(F\) from the same formula at \(\tilde F\). The term
\(\omega(h-\ell)\) cancels.
\begin{align*}
\Delta_\omega(\tilde F)-\Delta_\omega(F)&=(1-2\omega)\Bigl\{\E_{\tilde F}\bigl[(R-\ell)_+\bigr]-\E_F\bigl[(R-\ell)_+\bigr]\Bigr\}
\why{\eqref{eq:welfare-spreadpay} twice}\\
&=(1-2\omega)\int_\ell^{\bar r}\bigl[F(y)-\tilde F(y)\bigr]\,dy \why{Step~2}
\end{align*}
The integral is nonnegative. So the spread moves with the sign of \(1-2\omega\), weakly. It moves strictly exactly when
\(1-2\omega\ne0\) and the integral is positive. The integral is positive exactly when \(\E[(R-\ell)_+]\) changes
strictly.
\item \emph{The midpoint.} At \(\omega=1/2\) the coefficient \(1-2\omega\) is zero, so
\(\Delta_{1/2}=\tfrac12(h-\ell)\) at every incumbent distribution.
\end{enumerate}
\end{derivation}

\subsection{Why the direction depends on \texorpdfstring{\(\omega\)}{omega}}\label{sec:welfare-direction}

The proof gives the formula. The reason is also visible at a fixed incumbent value \(R\) with \(\ell<R<h\).
\begin{itemize}
\item Class \(H\) wins and pays \((1-\omega)R+\omega h\). The payment rises with \(R\) at the rate \(1-\omega\). The
target receives the runner-up's value, \(R\), with weight \(1-\omega\).
\item Class \(L\) loses to the incumbent, and the incumbent pays \((1-\omega)\ell+\omega R\). The payment rises with
\(R\) at the rate \(\omega\). The target receives the winner's value, \(R\), with weight \(\omega\).
\end{itemize}
A stronger incumbent raises \(R\). The class difference in the target's payment changes at the rate
\((1-\omega)-\omega=1-2\omega\).

Test the extremes.
\begin{itemize}
\item At \(\omega=0\), class \(L\)'s payment does not respond to \(R\), while class \(H\)'s payment follows \(R\) one
for one. The spread widens. This is the second-price rule, and it gives the point-mass spread of
Chapter~\ref{ch:auction}: \(\Delta_0=\E[(R-\ell)_+]=\int_\ell^{\bar r}[1-F(y)]\,dy\).
\item As \(\omega\) tends to one, class \(H\) pays almost \(h\), whatever \(R\) is. Class \(L\) pays almost
\(\max\{R,\ell\}\). The spread tends to \((h-\ell)-\E[(R-\ell)_+]\), which falls when the incumbent is stronger.
\item At \(\omega=1/2\) the two rates are equal, and the response vanishes.
\end{itemize}

\begin{watchout}
The weights give \emph{weak} directions of the spread. Suppose a stronger incumbent changes only the part of \(F\) below
\(\ell\). Then \(F=\tilde F\) on \([\ell,\bar r]\), the integral in Step~3 is zero, and \(\E[(R-\ell)_+]\) does not
change. The spread does not change at any \(\omega\), yet the challenger's profit falls strictly. A strict change in
the spread needs a strict change in the positive-part expectation, together with \(\omega\ne1/2\). The words ``a
stronger incumbent'' are not enough.
\end{watchout}

\begin{checkyourself}\label{q:welfare-strict}
Let \(\tilde F\le F\). (a)~Give a condition on \(\tilde F\) and \(F\) under which the spread \(\Delta_\omega\) does not
change at any \(\omega\), while the profit of a class-\(L\) challenger falls strictly. (b)~Give a condition under which
the spread falls strictly at \(\omega=\tfrac34\).
\end{checkyourself}

\begin{checkyourself}\label{q:welfare-pay0}
Set \(\omega=0\) in the point-mass case. Find the class difference in the target's payment at a fixed \(R\) with
\(\ell<R<h\), and compare it with \(\Gamma(R)=(R-\ell)_+\) of Chapter~\ref{ch:auction}. Why does the reserve not
appear in the comparison?
\end{checkyourself}

\subsection{The example}\label{sec:welfare-payexample}

In the example, the incumbent is uniform on \([0,r]\) with \(\ell<r<h\). Then \(\E[(R-\ell)_+]\) has a closed form.

\begin{derivation}[the positive part for a uniform incumbent]
\begin{align*}
\E\bigl[(R-\ell)_+\bigr]&=\int_\ell^r\Bigl(1-\frac yr\Bigr)\,dy \why{\eqref{eq:toolkit-pospart}, \(F(y)=y/r\) on \([0,r]\)}\\
&=(r-\ell)-\frac{r^2-\ell^2}{2r} \why{integrate}\\
&=\frac{(r-\ell)\,[2r-(r+\ell)]}{2r} \why{common denominator; \(r^2-\ell^2=(r-\ell)(r+\ell)\)}\\
&=\frac{(r-\ell)^2}{2r} \why{\(2r-(r+\ell)=r-\ell\)}
\end{align*}
\end{derivation}

This is the spread \(\DT(r)\) of the example at a reserve below \(\ell\). It is zero at \(r=\ell\), and it rises with
\(r\), because
\begin{align*}
\frac{d}{dr}\,\frac{(r-\ell)^2}{2r}&=\frac{2(r-\ell)\cdot2r-2(r-\ell)^2}{4r^2} \why{quotient rule}\\
&=\frac{(r-\ell)(r+\ell)}{2r^2}=\frac12\Bigl(1-\frac{\ell^2}{r^2}\Bigr)>0 \why{simplify; \(r>\ell\)}
\end{align*}
So \(\Delta_\omega(r)=\omega(h-\ell)+(1-2\omega)\DT(r)\). For two strengths \(\ell<r_0<r_1<h\),
\[
\Delta_\omega(r_1)-\Delta_\omega(r_0)=(1-2\omega)\bigl[\DT(r_1)-\DT(r_0)\bigr].
\]
The bracket is positive, so the change is zero exactly at \(\omega=1/2\). Chapter~\ref{app:numbers} gives its size for the
example's strengths.

The paper compares this result with Bulow and Klemperer (1996). They compare an auction with one more bidder against
an optimal negotiation, by expected revenue. The comparison here is different: it asks how competition moves the
spread.

Chapter~\ref{ch:ccc} needs a wider spread from the strong incumbent. Under \(\omega<1/2\), the payment stage moves the
spread in that direction. Under \(\omega>1/2\), it moves the spread the other way. The proposition concerns the
acquisition stage after every history. It states no equilibrium result.

% ====================================================================
\section{Reserves}\label{sec:welfare-reserves}

\subsection{What a reserve does}\label{sec:welfare-reservestory}

A reserve has two effects. It raises the payment that a winner must make. It also removes bidders whose value lies below
it. The second effect matters here. Take point-mass classes with \(\ell<p<h\), so that the reserve excludes the low
class. A low-class challenger can prepare, but it cannot bid.

Three events now differ. The paper defines them as follows.
\begin{itemize}
\item \emph{Preparation}: the challenger pays its preparation cost.
\item \emph{A sale}: at least one bidder submits an admissible bid.
\item \emph{Two admissible bidders}: the incumbent and a prepared challenger both meet the reserve.
\end{itemize}
When every value meets the reserve, a prepared challenger is always admissible, so preparation implies a sale. With a
reserve that excludes the low class, a low-class challenger can prepare and still produce no bid.

The probabilities show the difference. Write \(\Pr(\text{adm})\) for the probability that the challenger prepares and
meets the reserve. Given the class, preparation is independent of the values, so
\begin{align*}
\Pr(\text{adm})&=\sum_j\pi_j\,e_j\,\bigl[1-G_j(p^-)\bigr] \why{class weight, entry, chance to meet the reserve}\\
\Pr(\text{sale})&=\bigl[1-F(p)\bigr]+F(p)\,\Pr(\text{adm}) \why{the incumbent meets the reserve, or only the challenger does}\\
\Pr(\text{two})&=\bigl[1-F(p)\bigr]\,\Pr(\text{adm}) \why{both are admissible; independence}
\end{align*}
These three formulas are the guide's reading of the definitions. With a reserve that excludes the low class,
\(\Pr(\text{adm})=\pi\,e_H\), since a low-class challenger never meets the reserve.

\subsection{The spread with an excluded class}\label{sec:welfare-reservespread}

The reserve changes the spread. Chapter~\ref{ch:auction} derived the spread for any classes (\eqref{eq:auction-spread}):
\[
\DT(F)=p\,F(p)\,D(p^-)+\int_p^\infty\bigl[1-F(y)\bigr]D(y)\,dy .
\]
The first term is the \emph{reserve term}. The second is the \emph{slice term}. With a reserve below both classes,
\(D(p^-)=0\), and only the slice term is left. With a reserve that excludes the low class, \(D(p^-)=1\), and the reserve
term is \(pF(p)\).

\begin{intuition}[the reserve term]
Think of an incumbent that fails the reserve, which has probability \(F(p)\). An admissible challenger buys at the
reserve, and the target receives \(p\). An inadmissible challenger cannot bid, and the target stays unsold. The class
difference in the target's receipts is \(p\) times the difference \(D(p^-)\) in the chance to be admissible. This is the
reserve term. A weak incumbent fails the reserve often, so the term is large. A strong incumbent meets the reserve and
buys at \(p\), whatever the challenger does. The term shrinks.
\end{intuition}

To see the terms, compute the class objects for an incumbent uniform on \([0,r]\) and a point value \(\theta\). The paper
states the result as a kernel. The kernel gives the proceeds \(t(\theta)\), the profit \(g(\theta)\), and the chance \(W\)
that the challenger wins. Use the split \(t(\theta)+g(\theta)=\theta\) of \eqref{eq:auction-split} for \(\theta\ge p\),
and the profit formula \(g(\theta)=\int_p^\theta F(y)\,dy\) with \(F(y)=\min\{y/r,1\}\).

\begin{derivation}[the kernel for a uniform incumbent]
\emph{Case \(\theta<p\).} The challenger cannot bid, so \((t,g,W)=(t_0,0,0)\) with \(t_0=p(1-F(p))=p\,(1-p/r)_+\).

\emph{Case \(r\le p\le\theta\).} The incumbent never meets the reserve, so \(F(y)=1\) for \(y\ge p\).
\begin{align*}
g(\theta)&=\int_p^\theta1\,dy=\theta-p \why{\(F=1\) above \(p\)}\\
t(\theta)&=\theta-g(\theta)=p \why{\eqref{eq:auction-split}}\\
W&=F(\theta)=1 \why{\(\theta\ge p\ge r\)}
\end{align*}
\emph{Case \(p\le\theta\le r\).} Here \(F(y)=y/r\) on \([p,\theta]\).
\begin{align*}
g(\theta)&=\int_p^\theta\frac yr\,dy=\frac{\theta^2-p^2}{2r} \why{integrate}\\
t(\theta)&=\theta-\frac{\theta^2-p^2}{2r} \why{\eqref{eq:auction-split}}\\
W&=F(\theta)=\frac\theta r \why{\(F(y)=y/r\)}
\end{align*}
\emph{Case \(p\le r\le\theta\).} Split the integral at \(r\), where \(F\) reaches one.
\begin{align*}
g(\theta)&=\int_p^r\frac yr\,dy+\int_r^\theta1\,dy=\frac{r^2-p^2}{2r}+\theta-r \why{two pieces}\\
&=\theta-\frac r2-\frac{p^2}{2r} \why{\(\tfrac{r^2}{2r}-r=-\tfrac r2\)}\\
t(\theta)&=\frac r2+\frac{p^2}{2r} \why{\eqref{eq:auction-split}}\\
W&=1 \why{\(\theta\ge r\)}
\end{align*}
\end{derivation}

The cases agree where they meet. At \(\theta=r\), cases three and four both give \(g=(r^2-p^2)/(2r)\) and
\(t=r/2+p^2/(2r)\). At \(p=r\), case four gives \(g=\theta-r/2-r/2=\theta-p\) and \(t=r/2+r/2=p\), the values of case two.

Now take the point-mass classes with \(\ell<p<r<h\). Class \(L\) has the value \(\ell<p\), and class \(H\) has the value
\(h>r\). Then \(t_L=t_0=p-p^2/r\) and \(g_L=0\), and the fourth case gives \(t_H=r/2+p^2/(2r)\) and \(g_H=h-t_H\).
\begin{align*}
\DT&=t_H-t_L \why{definition}\\
&=\frac r2+\frac{p^2}{2r}-p+\frac{p^2}{r} \why{substitute}\\
&=\frac r2-p+\frac{3p^2}{2r} \why{\(\tfrac{p^2}{2r}+\tfrac{p^2}{r}=\tfrac{3p^2}{2r}\)}
\end{align*}
The formula \eqref{eq:auction-spread} gives the same spread in two terms. Here \(D(y)=\1\{\ell\le y<h\}\), \(D(p^-)=1\),
and \(F(p)=p/r\). The slice term has the integrand \(1-y/r\) on \([p,r]\) and zero above \(r\).
\begin{align*}
\DT&=p\cdot\frac pr+\int_p^r\Bigl(1-\frac yr\Bigr)dy \why{reserve term; slice term}\\
&=\frac{p^2}{r}+\frac{(r-p)^2}{2r} \why{the integral of Section~\ref{sec:welfare-payexample}, with \(p\) in place of \(\ell\)}
\end{align*}
The two forms agree: \(p^2/r+(r-p)^2/(2r)=(2p^2+r^2-2rp+p^2)/(2r)=r/2-p+3p^2/(2r)\). The two terms move in opposite
directions as the incumbent strengthens:
\begin{align*}
\frac{d}{dr}\,\frac{p^2}{r}&=-\frac{p^2}{r^2}<0 \why{reserve term}\\
\frac{d}{dr}\,\frac{(r-p)^2}{2r}&=\frac12\Bigl(1-\frac{p^2}{r^2}\Bigr)\ge0 \why{slice term, as in Section~\ref{sec:welfare-payexample}}
\end{align*}
The reserve term falls, since the incumbent meets the reserve more often. The slice term rises, as in
Chapter~\ref{ch:auction}. The slice term has zero slope at \(r=p\), so for a strength just above the reserve the first
effect wins and the spread falls. This agrees with part (c) of Proposition~1: equal admissibility fails here, and some
strengthening lowers the spread (Section~\ref{sec:auction-weakest}).

\begin{checkyourself}\label{q:welfare-min}
In the example above, find the strength \(r\) at which the spread \(\DT(r)=r/2-p+3p^2/(2r)\) is lowest, for \(r\ge p\),
and the lowest value, as a multiple of \(p\). Compare it with the spread at \(r=p\). Explain in words why the minimum is
above \(p\).
\end{checkyourself}

\subsection{The proposition}\label{sec:welfare-reservestatement}

The proof needs three facts about the class objects when class values may lie below the reserve. The paper's proposition
says to keep the class ranking and to allow class values below the reserve. The second half of the class ranking,
\(G_H(p^-)=G_L(p^-)=0\), cannot hold when values lie below the reserve. The guide reads the proposition as keeping the
first half, \(G_H\le G_L\). The proof uses only this half.

\begin{derivation}[facts used when class values may lie below the reserve]
Let \(G_H\le G_L\). Then \(D(y)\ge0\) for all \(y\). Also \(D(p^-)\ge0\), since the inequality holds at every left limit.
\begin{enumerate}[label=\textit{(F\arabic*)},leftmargin=*,itemsep=0.4em]
\item \emph{The spread is nonnegative.} The formula \eqref{eq:auction-spread} holds for any classes. Every factor is
nonnegative:
\[
\DT=\underbrace{p\,F(p)\,D(p^-)}_{\ge0}+\underbrace{\int_p^\infty[1-F(y)]\,D(y)\,dy}_{\ge0}\ge0 .
\]
\item \emph{Class \(H\) earns at least as much as class \(L\).} By \eqref{eq:auction-gross},
\(g_H-g_L=\int_p^\infty D(y)F(y)\,dy\ge0\). So \(B(\mu)\) is nondecreasing in \(\mu\).
\item \emph{Entry never lowers proceeds, and a strengthening never raises profit.} For every value \(\theta\), the target
receives at least \(p\,\1\{R\ge p\}\), so \(t_j\ge t_0\) and \(w_L\ge0\) (Section~\ref{sec:auction-outcomes}). And
Proposition~1(a) gives \(B_1\le B_0\) at every belief, for any classes.
\end{enumerate}
\end{derivation}

Proposition~1(b) is lost. Without equal admissibility, a strengthening can lower the spread, as the example above shows.
The proposition therefore asks for a large spread at both strengths.

\needspace{12\baselineskip}
\begin{paperresult}{The proposition on reserves (Appendix A.7 of the paper)}
Fix a reserve, and allow class values below it. Keep the class ranking, the costs, the slot and a strengthening
\(F_1\le F_0\), and compute each class object with the reserve. Suppose \(k<\rho\kappa\min\{\Delta_0,\Delta_1\}\).
\begin{enumerate}[label=(\alph*),leftmargin=2em]
\item If \(B_0(\pi)<c_H<B_1(M)\), entry is the same at the two strengths:
\(\Ent_0=\Ent_1=\rho+(1-\rho)(b+\pi\lambda)\).
\item If \(B_0(\pi)<c_H<B_0(M)\) and \(B_1(M)<c_H\), entry falls strictly:
\(\Ent_0=\rho+(1-\rho)(b+\pi\lambda)>\rho=\Ent_1\).
\end{enumerate}
Each economy has a unique outcome under every pair of tie rules.
\end{paperresult}

Read the statement line by line.
\begin{itemize}
\item \emph{The spread condition.} It replaces the window (H2) of Proposition~2. Since \(k>0\), it gives \(\Delta_0>0\)
and \(\Delta_1>0\). It names both strengths, because the direction of the spread is no longer known.
\item \emph{Part (a).} A purchase recruits at the strong incumbent, so it recruits at the weak incumbent as well. Entry
is the same: \(\rho\), plus the chance \(1-\rho\) of a high cost times the chance \(b+\pi\lambda\) of a purchase.
\item \emph{Part (b).} A purchase recruits at the weak incumbent but not at the strong incumbent. Entry falls to the
floor.
\item \emph{Tie rules.} All four pairs give the same outcome. Every comparison in the proof is strict.
\item \emph{The slot.} The statement keeps the slot of the model and gives no separate domain for \(\beta\). The proof
uses only the posterior bounds after nonzero orders.
\end{itemize}

\subsection{The proof}\label{sec:welfare-reserveproof}

\begin{derivation}[proof of the proposition on reserves]
\begin{enumerate}[label=\textit{Step \arabic*.},leftmargin=*,itemsep=0.6em]
\item \emph{The orders are full at both strengths.} Fix a strength \(i\in\{0,1\}\) and any candidate. The trading lemma
of Chapter~\ref{ch:trading} (Lemma~A.3) uses only Bayes' rule for the posteriors, competitive prices, truthful bids,
entry in \([\rho,1]\), and \(\DT\ge0\). All hold with the reserve, by (F1). The trading bound gives
\begin{align*}
K_H&\ge\rho\,\kappa\,\Delta_i \why{trading bound \eqref{eq:trading-bound}}\\
&>k \why{hypothesis, \(k<\rho\kappa\min\{\Delta_0,\Delta_1\}\)}
\end{align*}
The same holds for \(K_L\). A wrong-signed order earns at most \(-k\). So each class has one best order, the correct one.
The orders are full, and the posteriors are \((m,\pi,M)\). The bounds hold for every candidate. The spread is positive,
since \(\Delta_i>k/(\rho\kappa)>0\).

\item \emph{The three prices differ, so the beliefs are the posteriors.} Suppose two orders \(x\) and \(x'\) had one
price. Entry depends on the order only through the price, so \(e(x)=e(x')=\bar e\ge\rho>0\). The price formula gives
\(0=P(x)-P(x')=\bar e\,\Delta_i\,[\mu(x)-\mu(x')]\). Then \(\mu(x)=\mu(x')\), which is false, since \(m<\pi<M\). So each
price has one order, and the belief at the price is the posterior of its order.

\item \emph{Entry at the weak incumbent.} Both parts have \(B_0(\pi)<c_H\). Use (F2) and (F3).
\begin{align*}
B_0(m)&\le B_0(\pi) \why{(F2), \(m<\pi\)}\\
&<c_H \why{hypothesis}
\end{align*}
After a sale and after no order, gross profit is below \(c_H\), so \(e=\rho\). After a purchase the belief is \(M\). In
part (a), \(B_0(M)\ge B_1(M)>c_H\) by (F3) and the hypothesis. In part (b), \(B_0(M)>c_H\) by the hypothesis. Either way,
a purchase recruits, and \(e(1)=1\). Entry at \(F_0\) is \((\rho,\rho,1)\).

\item \emph{Entry at the strong incumbent.} In part (a), \(B_1(\pi)\le B_0(\pi)<c_H<B_1(M)\) by (F3). The argument of
Step~3 then gives entry \((\rho,\rho,1)\) at \(F_1\). In part (b), \(B_1(M)<c_H\). For every posterior \(\mu\le M\),
\begin{align*}
B_1(\mu)&\le B_1(M) \why{(F2); \(\mu\le M\)}\\
&<c_H \why{hypothesis}
\end{align*}
No price recruits, and \(e\equiv\rho\) at \(F_1\).

\item \emph{Total entry.} Entry \((\rho,\rho,1)\) means \(e(x)=\rho+(1-\rho)\1\{x=1\}\). Then
\begin{align*}
\Ent&=\sum_x\Pr(X=x)\,e(x) \why{total probability over orders}\\
&=\rho+(1-\rho)\Pr(X=1) \why{\(\sum_x\Pr(X=x)=1\)}\\
&=\rho+(1-\rho)(b+\pi\lambda) \why{\(\Pr(X=1)=\pi(b+\lambda)+(1-\pi)b\)}
\end{align*}
This is \eqref{eq:ccc-entry}. Entry \(e\equiv\rho\) gives \(\Ent=\rho\). Part (a) has the first case at both strengths.
Part (b) has the first case at \(F_0\) and the second at \(F_1\). Since \((1-\rho)(b+\pi\lambda)>0\), part (b) shows a
strict fall.

\item \emph{Equilibrium and tie rules.} Each tuple has full orders, the posteriors \((m,\pi,M)\), this entry, the prices
\(P_i(x)=t_{0,i}+e_i(x)\,w_i(\mu(x))\), the beliefs of Step~2 and truthful bids. Bayes' rule holds by the order
probabilities, and the high-cost choice is strict at every price. The prices are competitive by the price formula, the
orders are strict best orders by Step~1, and the bids are truthful. Every comparison is strict, so no tie rule is used.
Steps~1 to~4 hold for every candidate, so each tuple is the only equilibrium, under every pair of tie rules.
\end{enumerate}
\end{derivation}

\begin{watchout}
Strict comparisons carry the tie-rule statement. At a boundary reserve where gross profit at \(M\) equals \(c_H\), a
purchase leaves a high-cost challenger indifferent. The challenger's tie rule then matters: the challenger prepares under
the default rule, and it may prepare with any probability under the other. The proposition excludes this case with strict
inequalities. The proposition also compares equilibria. Each economy is solved again, and a deviation faces the prices
and entry of its own economy.
\end{watchout}

\subsection{Why entry stays flat}\label{sec:welfare-flat}

Compare part (a) with Proposition~2. Proposition~2 needs \(\Delta_0<k\): the weak incumbent leaves a spread too small to
pay for an informed order. A strengthening then switches information on, and a purchase recruits at the strong incumbent
only. Entry rises.

A reserve that excludes the low class can break the left half of (H2). The reserve term gives
\begin{align*}
\Delta_0&\ge p\,F_0(p)\,D(p^-) \why{the slice term in \eqref{eq:auction-spread} is nonnegative}\\
&=p\,F_0(p) \why{\(D(p^-)=1\) with the low class excluded}
\end{align*}
When the weak incumbent often fails the reserve, \(pF_0(p)\) is large. The spread already pays for informed trade, and the
price speaks at the weak incumbent. A purchase recruits there as well, as part (a) says. Competition has no information to
switch on, and entry is flat.

Only a fall in gross profit can then change entry. This is part (b): a strong incumbent can push gross profit at the top
posterior below \(c_H\). In Proposition~2, a strengthening raised entry from \(F_0\) to \(F_1\). In the two cases of the
reserve proposition, entry never rises with incumbent strength.

\begin{checkyourself}\label{q:welfare-h2}
Take the point-mass classes with \(\ell<p<h\). Show that \(\Delta_0\ge pF_0(p)\). Suppose \(pF_0(p)>k\). Which inequality
of (H2) fails, and what does the investor do against the weak incumbent?
\end{checkyourself}

\subsection{Reserves in the example}\label{sec:welfare-reserveexample}

The paper reports three reserves in the example. Chapter~\ref{app:numbers} gives the numbers.
\begin{itemize}
\item \emph{A reserve below both challenger values.} Every prepared challenger is admissible. This is the setting of
Chapter~\ref{ch:ccc}, and entry first rises with incumbent strength and then falls.
\item \emph{A reserve that excludes the low class and lies below the weak incumbent's top.} Part (a) applies. Full orders
are the unique outcome at both strengths, only a purchase recruits, and entry is the same at both. This reserve raises
expected proceeds at both strengths, relative to the reserve below both values.
\item \emph{A reserve above the weak incumbent's top and below \(h\).} Part (b) applies in a window. Entry is the floor
at the strong incumbent, and it is higher at the weak incumbent. Proceeds exceed their levels at the reserve below both
values.
\end{itemize}
The paper reports that, in the example, strict deterrence with a reserve that excludes the low class needs the reserve to
exceed the weak incumbent's top. For a reserve between \(\ell\) and that top, gross profit at the top posterior exceeds
the high cost at both strengths, and entry is flat.

The window can be written in symbols. For the third reserve, the low class earns nothing, so \(B_i(\mu)=\mu\,g_{H,i}\)
and \(B_i(M)=M\,g_{H,i}\). At the strong incumbent, \(p\le r_1\le h\), and the fourth case of the kernel gives
\(g_{H,1}=h-r_1/2-p^2/(2r_1)\). At the weak incumbent, \(r_0\le p\le h\), and the second case gives \(g_{H,0}=h-p\).

\begin{derivation}[the deterrence window]
\emph{Strong incumbent.} Gross profit at the top posterior is below the high cost when
\begin{align*}
M\Bigl(h-\frac{r_1}2-\frac{p^2}{2r_1}\Bigr)&<c_H \why{part (b), \(B_1(M)<c_H\)}\\
\frac{p^2}{2r_1}&>h-\frac{c_H}M-\frac{r_1}2 \why{divide by \(M>0\); rearrange}\\
p^2&>2r_1\Bigl(h-\frac{c_H}M\Bigr)-r_1^2 \why{multiply by \(2r_1>0\)}
\end{align*}
The right side defines \(p_H^2\). When it is positive, this is \(p>p_H\) with
\(p_H=\sqrt{2r_1(h-c_H/M)-r_1^2}\).

\emph{Weak incumbent.} Gross profit at the top posterior exceeds the high cost when
\begin{align*}
M\,(h-p)&>c_H \why{part (b), \(B_0(M)>c_H\)}\\
p&<h-\frac{c_H}M \why{divide by \(M>0\); rearrange}
\end{align*}
\end{derivation}

So part (b) holds in the example for \(p_H<p<h-c_H/M\), together with the prior condition \(B_0(\pi)<c_H\) and with
\(p>r_0\). A higher reserve lowers gross profit at both strengths. The paper checks the three margins at its reserve and
reports that \(p_H\) exceeds \(r_0\).

% ====================================================================
\section{The three results together}\label{sec:welfare-together}

The three results share one structure. The price carries information through the spread. A term of the sale changes the
spread, or changes who sees the price, and entry follows.
\begin{itemize}
\item Access adds the right high-cost challengers and no others. It raises target proceeds and net surplus. The proof
counts the added preparations and weights each by its gain.
\item The payment rule fixes how the spread responds to competition. The response has the sign of \(1-2\omega\). A
strict move also needs the positive-part expectation to change strictly.
\item A reserve that excludes the low class adds a reserve term to the spread. A stronger incumbent shrinks this term.
The weak incumbent already leaves a wide spread, so entry is flat unless a stronger incumbent deters preparation.
\end{itemize}

\begin{chaptersummary}
\begin{itemize}
\item When the challenger sees the price, high-cost challengers prepare after a purchase and nowhere else. The orders and
the trading costs are the same with and without access.
\item Each added preparation raises target proceeds by \(w_1(M)\) on average. It raises net surplus by
\(B_1(M)-c_H+pF_1(p)\), which is positive since \(B_1(M)>c_H\).
\item The informed investor's gross profit is a transfer from liquidity traders. The shareholders still gain when one order
unit is below the ratio of the proceeds gain to this profit. A constant dividend shows that the information in the price,
and not its level, draws entry.
\item Outside the recruitment region, access removes preparation after a sale. Proceeds fall, and the sign of the surplus
change depends on \(B_1(m)-c_H+pF_1(p)\).
\item Under Nash bargaining with the seller weight \(\omega\), the spread is
\(\omega(h-\ell)+(1-2\omega)\E[(R-\ell)_+]\). A stronger incumbent never raises the challenger's profit. It moves the
spread weakly in the direction of \(1-2\omega\). A strict move needs a strict change in \(\E[(R-\ell)_+]\) and
\(\omega\ne\tfrac12\).
\item With a reserve that excludes the low class, the spread has a reserve term \(pF(p)\), which falls with incumbent
strength. If informed trade is full at both strengths and only a purchase recruits at each, entry is flat. Strict
deterrence needs a stronger incumbent to push \(B_1(M)\) below \(c_H\).
\end{itemize}
\end{chaptersummary}

% ====================================================================
\answersection

\answer{\ref{q:welfare-sign}}{%
Outside the recruitment region, access removes high-cost preparation after a sale. Entry with access is \((\rho,1,1)\),
and entry without access is one after every order. The class entries with access are
\(e_H^1=1-(1-\rho)b\) and \(e_L^1=1-(1-\rho)(b+\lambda)\) (Section~\ref{sec:welfare-outside}). So the entry changes are
\(-(1-\rho)b\) and \(-(1-\rho)(b+\lambda)\). Use \eqref{eq:welfare-W} in both economies. The terms \(v_0\), and the
terms with \(\rho\), cancel in the difference:
\begin{align*}
\mathcal W^1-\mathcal W^0&=\sum_j\pi_j\,(a_j-c_H)\,(e_j^1-e_j^0) \why{\eqref{eq:welfare-W} twice; collect \(e_j\)}\\
&=-(1-\rho)\bigl\{\pi\,b\,(a_H-c_H)+(1-\pi)(b+\lambda)(a_L-c_H)\bigr\} \why{the entry changes}
\end{align*}
A sale has probability \(\Pr(X=-1)=\pi b+(1-\pi)(b+\lambda)=b+(1-\pi)\lambda\). Bayes' rule gives the weights
\(\pi b=(b+(1-\pi)\lambda)\,m\) and \((1-\pi)(b+\lambda)=(b+(1-\pi)\lambda)(1-m)\). Then
\begin{align*}
\mathcal W^1-\mathcal W^0&=-(1-\rho)\bigl[b+(1-\pi)\lambda\bigr]\bigl\{m\,a_H+(1-m)\,a_L-c_H\bigr\}
\why{the weights; \(m+(1-m)=1\)}\\
&=-(1-\rho)\bigl[b+(1-\pi)\lambda\bigr]\bigl[B_1(m)+pF_1(p)-c_H\bigr] \why{\eqref{eq:welfare-aj}}
\end{align*}
The prefactor is negative. So access raises surplus exactly when \(B_1(m)+pF_1(p)<c_H\), and it lowers surplus exactly
when \(B_1(m)+pF_1(p)>c_H\). We have \(B_1(m)<c_H\), so the first part is negative. The sale term \(pF_1(p)\) is
nonnegative, and it is large when the incumbent often fails the reserve. A removed preparation then costs the economy a
sale that it would not otherwise have. The challenger does not count this sale: it counts only its own profit
\(B_1(m)\). So a removed preparation can be worth its cost to the economy and not to the challenger.}

\answer{\ref{q:welfare-strict}}{%
(a) The spread changes by \((1-2\omega)\int_\ell^{\bar r}[F-\tilde F]\,dy\). It is unchanged at every \(\omega\) when
\(\tilde F=F\) on \([\ell,\bar r]\). Take also \(\tilde F<F\) somewhere in \((0,\ell)\). The profit of a class-\(L\)
challenger changes by
\begin{align*}
g_{\ell,\omega}(F)-g_{\ell,\omega}(\tilde F)&=(1-\omega)\int_0^\ell\bigl[F(y)-\tilde F(y)\bigr]\,dy
\why{Step~1 of the proof, \(\theta=\ell\)}\\
&>0 \why{\(1-\omega>0\); the integrand is continuous, nonnegative, and positive somewhere}
\end{align*}
The profit falls strictly, and the spread does not change. The stronger incumbent moves only values below \(\ell\), and
they never enter the spread.

(b) At \(\omega=\tfrac34\), \(1-2\omega=-\tfrac12<0\). The spread changes by
\(-\tfrac12\int_\ell^{\bar r}[F-\tilde F]\,dy\). It falls strictly exactly when the integral is positive. The integrand is
continuous and nonnegative, so this holds exactly when \(\tilde F<F\) at some point of \((\ell,\bar r)\). The seller keeps
three quarters of the surplus above the fallback. A stronger incumbent then raises the payment for class \(L\) by more
than the payment for class \(H\), and the class difference falls.}

\answer{\ref{q:welfare-pay0}}{%
At \(\omega=0\) the transfer is \(T_0(R,\theta)=\min\{R,\theta\}\). For \(\ell<R<h\),
\begin{align*}
T_0(R,h)-T_0(R,\ell)&=\min\{R,h\}-\min\{R,\ell\} \why{\(\omega=0\) in \eqref{eq:welfare-transfer}}\\
&=R-\ell \why{\(R<h\) and \(R>\ell\)}
\end{align*}
Chapter~\ref{ch:auction} found \(\Gamma(R)=\int_p^RD(y)\,dy=R-\ell\) for \(\ell\le R<h\) in the point-mass case with
\(p<\ell\). The two agree. The setting here has a zero reserve. In the main model, a reserve \(p<\ell\) enters \(\Gamma\)
only as the lower limit of the integral. Since \(D(y)=\1\{\ell\le y<h\}\) is zero on \([p,\ell)\), every \(p<\ell\) gives the
same \(\Gamma\), including the limit \(p\to0\). For \(R\le\ell\) both give zero, as the formula \((R-\ell)_+\) shows.}

\answer{\ref{q:welfare-min}}{%
Differentiate the spread with respect to \(r\).
\begin{align*}
\frac{d\DT}{dr}&=\frac12-\frac{3p^2}{2r^2} \why{derivative of \(\tfrac r2-p+\tfrac{3p^2}{2r}\)}\\
&=\frac{r^2-3p^2}{2r^2} \why{common denominator}
\end{align*}
The derivative is negative for \(p\le r<\sqrt3\,p\), zero at \(r=\sqrt3\,p\), and positive above. So \(\DT\) falls and then
rises, and the minimum is at \(r=\sqrt3\,p\). Its value is
\begin{align*}
\DT(\sqrt3\,p)&=\frac{\sqrt3\,p}2-p+\frac{3p^2}{2\sqrt3\,p} \why{substitute \(r=\sqrt3\,p\)}\\
&=\frac{\sqrt3}2\,p-p+\frac{\sqrt3}2\,p \why{\(3/(2\sqrt3)=\sqrt3/2\)}\\
&=(\sqrt3-1)\,p \why{add}
\end{align*}
At \(r=p\) the spread is \(p/2-p+3p/2=p\). Since \(\sqrt3-1<1\), the minimum lies below this value. The minimum lies above
\(p\) because the two terms of the spread balance there. At \(r=p\), the incumbent fails the reserve for sure, and the
reserve term is largest. The slice term is zero and has zero slope. As \(r\) rises, the reserve term falls at the rate
\(p^2/r^2\), and the slice term rises at the rate \(\tfrac12(1-p^2/r^2)\), which starts at zero. The first rate exceeds the
second until \(r=\sqrt3\,p\).}

\answer{\ref{q:welfare-h2}}{%
With point masses and \(\ell<p<h\), \(G_L(p^-)=1\) and \(G_H(p^-)=0\), so \(D(p^-)=1\). By \eqref{eq:auction-spread},
\begin{align*}
\Delta_0&=p\,F_0(p)\cdot1+\int_p^\infty\bigl[1-F_0(y)\bigr]D(y)\,dy \why{\eqref{eq:auction-spread}, \(D(p^-)=1\)}\\
&\ge p\,F_0(p) \why{the integrand is nonnegative: \(F_0\le1\), \(D\ge0\)}
\end{align*}
Suppose \(pF_0(p)>k\). Then \(\Delta_0>k\), so the left inequality \(\Delta_0<k\) of (H2) fails. The upper bound of the
trading return no longer rules out informed orders at the weak incumbent. If also \(k<\rho\kappa\Delta_0\), the trading
bound gives \(K_H\ge\rho\kappa\Delta_0>k\) and \(K_L\ge\rho\kappa\Delta_0>k\). Then class \(H\) buys and class \(L\) sells at
the weak incumbent as well, the price speaks there, and a purchase recruits when \(B_0(M)>c_H\). Informed trade does not
need a stronger incumbent.}
